\documentclass[12pt,letterpaper]{article}
\usepackage{graphicx,textcomp}
\usepackage{natbib}
\usepackage{setspace}
\usepackage{fullpage}
\usepackage{color}
\usepackage[reqno]{amsmath}
\usepackage{amsthm}
\usepackage{fancyvrb}
\usepackage{amssymb,enumerate}
\usepackage[all]{xy}
\usepackage{endnotes}
\usepackage{lscape}
\newtheorem{com}{Comment}
\usepackage{float}
\usepackage{hyperref}
\newtheorem{lem} {Lemma}
\newtheorem{prop}{Proposition}
\newtheorem{thm}{Theorem}
\newtheorem{defn}{Definition}
\newtheorem{cor}{Corollary}
\newtheorem{obs}{Observation}
\usepackage[compact]{titlesec}
\usepackage{dcolumn}
\usepackage{tikz}
\usetikzlibrary{arrows}
\usepackage{multirow}
\usepackage{xcolor}
\newcolumntype{.}{D{.}{.}{-1}}
\newcolumntype{d}[1]{D{.}{.}{#1}}
\definecolor{light-gray}{gray}{0.65}
\usepackage{url}
\usepackage{listings}
\usepackage{color}

% 改善段落间距和对齐
\setlength{\parindent}{0pt}
\setlength{\parskip}{6pt plus 2pt minus 1pt}
\raggedright

\definecolor{codegreen}{rgb}{0,0.6,0}
\definecolor{codegray}{rgb}{0.5,0.5,0.5}
\definecolor{codepurple}{rgb}{0.58,0,0.82}
\definecolor{backcolour}{rgb}{0.95,0.95,0.92}

\lstdefinestyle{mystyle}{
	backgroundcolor=\color{backcolour},   
	commentstyle=\color{codegreen},
	keywordstyle=\color{magenta},
	numberstyle=\tiny\color{codegray},
	stringstyle=\color{codepurple},
	basicstyle=\footnotesize,
	breakatwhitespace=false,         
	breaklines=true,                 
	captionpos=b,                    
	keepspaces=true,                 
	numbers=left,                    
	numbersep=5pt,                  
	showspaces=false,                
	showstringspaces=false,
	showtabs=false,                  
	tabsize=2
}
\lstset{style=mystyle}
\newcommand{\Sref}[1]{Section~\ref{#1}}
\newtheorem{hyp}{Hypothesis}

\title{Problem Set 4}
\date{\today}
\author{Hanyu Li (ID: 25346841)}

\begin{document}
	\maketitle
	
	\section*{Instructions}
	\begin{itemize}
		\item Please show your work! You may lose points by simply writing in the answer. If the problem requires you to execute commands in \texttt{R}, please include the code you used to get your answers. Please also include the \texttt{.R} file that contains your code. If you are not sure if work needs to be shown for a particular problem, please ask.
		\item Your homework should be submitted electronically on GitHub.
		\item This problem set is due before 23:59 on Thursday November 27, 2025. No late assignments will be accepted.
	\end{itemize}
	
	\vspace{.5cm}
	
	\section*{Question 1: Economics}
	\vspace{.25cm}
	
	\noindent In this question, use the \texttt{prestige} dataset in the \texttt{car} library. First, run the following commands:
	
	\begin{verbatim}
		install.packages(car)
		library(car)
		data(Prestige)
		help(Prestige)
	\end{verbatim}
	
	\noindent We would like to study whether individuals with higher levels of income have more prestigious jobs. Moreover, we would like to study whether professionals have more prestigious jobs than blue and white collar workers.
	
	\newpage
	
	\begin{enumerate}
		\item [(a)] Create a new variable \texttt{professional} by recoding the variable \texttt{type} so that professionals are coded as $1$, and blue and white collar workers are coded as $0$ (Hint: \texttt{ifelse}).
		
		\vspace{.3cm}
		\lstinputlisting[language=R, firstline=33, lastline=42]{PS04_HL.R}
		
		\vspace{.3cm}
		Check if we add variable \texttt{professional} successfully:
		
		\begin{verbatim}
			                    education income women prestige census type professional
			gov.administrators      13.11  12351 11.16     68.8   1113 prof    1
			general.managers        12.26  25879  4.02     69.1   1130 prof    1
			accountants             12.77   9271 15.70     63.4   1171 prof    1
			purchasing.officers     11.42   8865  9.11     56.8   1175 prof    1
			chemists                14.62   8403 11.68     73.5   2111 prof    1
			physicists              15.64  11030  5.13     77.6   2113 prof    1
		\end{verbatim}
		
		\item [(b)] Run a linear model with \texttt{prestige} as an outcome and \texttt{income}, \texttt{professional}, and the interaction of the two as predictors (Note: this is a continuous $\times$ dummy interaction).
		
		\vspace{.3cm}
		To build an interaction model, since we have professional as a dummy variable, we assume:
		
		\[
		prestige_i = \beta_0 + \beta_1\,income_i + \beta_2\,professional_i + \beta_3\,(income_i \times professional_i) + \varepsilon_i.
		\]
		
		To test whether the interaction term is statistically significant, we consider the following hypotheses:
		
		\[
		H_0:\ \beta_3 = 0 \quad \text{versus} \quad H_a:\ \beta_3 \neq 0.
		\]
		
		Here I set unprofessional respondents as the baseline group and the interaction model is built as below:
		
		\lstinputlisting[language=R, firstline=44, lastline=52]{PS04_HL.R}
		
		After running the regression we have this result:
		
		\begin{table}[H] 
			\centering
			\label{tab:interaction}
			\input{interact_reg.tex}
		\end{table}
		
		Here the p-value for the interaction term $(income \times professional)$ is less than 0.001, which means we have statistical evidence to reject $H_0$.
		
		We can thus interpret that the effect of income on occupational prestige differs for professionals and non-professionals. Specifically, for every one unit increase in income, the marginal increase in prestige scores is 0.00233 points smaller for professionals than for non-professionals.
		
		The model explains approximately 78.7\% of the variation in occupational prestige scores, indicating that the predictors included in the model account for a large share of the variation in the outcome.
		
		\item [(c)] Write the prediction equation based on the result.
		
		\vspace{.3cm}
		\[
		\widehat{prestige}_i = 21.14226 + 0.00317\,income_i + 37.78128\,professional_i - 0.00233\,(income_i \times professional_i).
		\]
		
		\item [(d)] Interpret the coefficient for \texttt{income}.
		
		\vspace{.3cm}
		The coefficient 0.00317 indicates: for non-professionals (the baseline \texttt{professional}=0), one unit increase in income is associated with an average increase of 0.00317 points in the Pineo-Porter prestige score.
		
		\item [(e)] Interpret the coefficient for \texttt{professional}.
		
		\vspace{.3cm}
		The coefficient 37.78128 indicates: when income equals 0, professional occupations have, on average, prestige scores that are 37.78 points higher than non-professional occupations.
		
		\item [(f)] What is the effect of a \$1,000 increase in income on prestige score for professional occupations? In other words, we are interested in the marginal effect of income when the variable \texttt{professional} takes the value of $1$. Calculate the change in $\hat{y}$ associated with a \$1,000 increase in income based on your answer for (c).
		
		\vspace{.3cm}
		My whole thinking process is calculating when $professional_i = 1$, the change of prestige score for a \$1,000 increase in income:
		
		\[
		\widehat{prestige}_i = 21.14226 + 0.00317 \, income_i + 37.78128 \, professional_i - 0.00233 \, (income_i \times professional_i)
		\]
		
		For the above model, we substitute $professional_i = 1$ and simplify:
		
		\[
		\widehat{prestige}_i = 21.14226 + 0.00317 \, income_i + 37.78128 - 0.00233 \, income_i
		\]
		
		\[
		\widehat{prestige}_i = 58.92354 + 0.00084 \, income_i
		\]
		
		To compute the change in the predicted prestige score associated with a \$1,000 increase in income, we compute the difference:
		
		\[
		\Delta prestige = \widehat{prestige}(income_i + 1000) - \widehat{prestige}(income_i)
		\]
		
		\[
		\Delta prestige = 0.00084 \times 1000 = 0.84
		\]
		
		Therefore, for professionals, a \$1,000 increase in income is associated with an increase of about 0.84 points in the predicted prestige score.
		
		\item [(g)] What is the effect of changing one's occupations from non-professional to professional when her income is \$6,000? We are interested in the marginal effect of professional jobs when the variable \texttt{income} takes the value of $6,000$. Calculate the change in $\hat{y}$ based on your answer for (c).
		
		\vspace{.3cm}
		My whole thinking process is calculating the difference of prestige scores between $professional_i = 1$ and $professional_i = 0$, when \texttt{income} takes the value of $6,000$.
		
		\[
		\Delta \widehat{prestige} = \widehat{prestige}_{\text{pro}} - \widehat{prestige}_{\text{non-pro}}
		\]
		
		\[
		\widehat{prestige}_{\text{pro}} = 21.14226 + 0.00317 \cdot 6000 + 37.78128 \cdot 1 - 0.00233 \cdot (6000 \cdot 1) = 63.96354
		\]
		
		\[
		\widehat{prestige}_{\text{non-pro}} = 21.14226 + 0.00317 \cdot 6000 + 37.78128 \cdot 0 - 0.00233 \cdot (6000 \cdot 0) = 40.16226
		\]
		
		\[
		\begin{aligned}
			\Delta \widehat{prestige} = 63.96354 - 40.16226 = 23.80128
		\end{aligned}
		\]
		
		Thus, when someone with an income of \$6,000 changes from non-professional to professional occupation, there will be an approximate 23.80 points increase in predicated prestige score.
		
	\end{enumerate}
	
	\newpage
	
	\section*{Question 2: Political Science}
	\vspace{.25cm}
	
	\noindent Researchers are interested in learning the effect of all of those yard signs on voting preferences.\footnote{Donald P. Green, Jonathan	S. Krasno, Alexander Coppock, Benjamin D. Farrer,	Brandon Lenoir, Joshua N. Zingher. 2016. ``The effects of lawn signs on vote outcomes: Results from four randomized field experiments.'' Electoral Studies 41: 143-150. } Working with a campaign in Fairfax County, Virginia, 131 precincts were randomly divided into a treatment and control group. In 30 precincts, signs were posted around the precinct that read, ``For Sale: Terry McAuliffe. Don't Sellout Virgina on November 5.''
	
	Below is the result of a regression with two variables and a constant. The dependent variable is the proportion of the vote that went to McAuliff's opponent Ken Cuccinelli. The first variable indicates whether a precinct was randomly assigned to have the sign against McAuliffe posted. The second variable indicates a precinct that was adjacent to a precinct in the treatment group (since people in those precincts might be exposed to the signs).
	
	\vspace{.5cm}
	\begin{table}[!htbp]
		\centering 
		\textbf{Impact of lawn signs on vote share}\\
		\begin{tabular}{@{\extracolsep{5pt}}lccc} 
			\\[-1.8ex] 
			\hline \\[-1.8ex]
			Precinct assigned lawn signs  (n=30)  & 0.042\\
			& (0.016) \\
			Precinct adjacent to lawn signs (n=76) & 0.042 \\
			&  (0.013) \\
			Constant  & 0.302\\
			& (0.011)
			\\
			\hline \\
		\end{tabular}\\
		\footnotesize{\textit{Notes:} $R^2$=0.094, N=131}
	\end{table}
	
	\vspace{.5cm}
	
	\begin{enumerate}
		\item [(a)] Use the results from a linear regression to determine whether having these yard signs in a precinct affects vote share (e.g., conduct a hypothesis test with $\alpha = .05$).
		
		\vspace{.3cm}
		From the regression outcome we can write the prediction equation as:
		
		\[
		\widehat{vote\_share}_i = 0.302 + 0.042\, precinctAssigned_i + 0.042\, precinctAdjacent_i.
		\]
		
		We have the following estimated coefficients:
		
		\[
		\beta_0 \text{ (intercept)} = 0.302,\quad \beta_1 \text{ (coefficient of precinct assigned lawn signs)} = 0.042,\quad \beta_2 \text{ (coefficient of precinct adjacent to lawn signs)} = 0.042.
		\]
		
		To conduct a hypothesis test, we first list our hypothesis:
		
		\[
		H_0: \beta_1 = 0 \qquad\text{vs.}\qquad H_a: \beta_1 \neq 0.
		\]
		
		Then, we calculate our test statistic. Since the population variance is unknown, the $t$ distribution is used here. According to the formula below, we calculate the test statistic as $\textit{t-value}_1 = 2.625$:
		
		\[
		t = \frac{\beta_1 - 0}{SE(\beta_1)}.
		\]
		
		For the calculation of the critical value, since this is a two–tailed test at the 0.05 significance level, and the model contains 1 intercept and 2 independent variables (3 estimators in total), the degrees of freedom are:
		
		\[
		df = N - k = 131 - 3 = 128.
		\]
		
		Therefore, the critical value is calculated as follows, and the result is:
		\[
		t_{\text{critical}} = 1.978671.
		\]
		
		\lstinputlisting[language=R, firstline=64, lastline=74]{PS04_HL.R}
		
		Thirdly, we calculate the p-value and the result is:
		\[
		p\text{-value}_1 = 0.00972.
		\]
		
		\lstinputlisting[language=R, firstline=76, lastline=78]{PS04_HL.R}
		
		In summary, since $p\text{-value}_1 = 0.00972$ is lower than the decided threshold 0.05, and the test statistic $\textit{t-value}_1 = 2.625$ exceeds the two–tailed 5\% critical value of 1.978671, we have sufficient evidence to reject $H_0$ and accept $H_a$, concluding that having yard signs in a precinct does affect the vote share.
		
		\item [(b)] Use the results to determine whether being next to precincts with these yard signs affects vote share (e.g., conduct a hypothesis test with $\alpha = .05$).
		
		As the same steps we have gone through in question (a), firstly we have our hypothesis:
		
		\[
		H_0: \beta_2 = 0 \qquad\text{vs.}\qquad H_a: \beta_2 \neq 0.
		\]
		
		Then, we calculate the test statistic, $\textit{t-value}_2 = 3.230769$:
		
		\[
		t = \frac{\beta_2 - 0}{SE(\beta_2)}.
		\]
		
		Since the hypothesis test for $\beta_2$ shares the same degrees of freedom and significance level, the critical value is still:
		\[
		t_{\text{critical}} = 1.978671.
		\]
		
		\lstinputlisting[language=R, firstline=80, lastline=88]{PS04_HL.R}
		
		Thirdly, we calculate the p-value and the result is:
		\[
		p\text{-value}_2 = 0.001569462.
		\]
		
		\lstinputlisting[language=R, firstline=90, lastline=92]{PS04_HL.R}
		
		To sum, since $p\text{-value}_2 = 0.001569462$ is lower than the decided threshold 0.05, and the test statistic $\textit{t-value}_2 = 3.230769$ exceeds the two–tailed 5\% critical value of 1.978671, we have sufficient evidence to reject $H_0$ and accept $H_a$, concluding that being next to precincts with these yard signs also affects the vote share.
		
		\item [(c)] Interpret the coefficient for the constant term substantively.
		
		I interpret the constant coefficient (0.302) as follows: For precincts that were neither assigned lawn signs against McAuliffe nor adjacent to precincts with such signs, the predicted vote share going to McAuliffe's opponent, Ken Cuccinelli, is approximately 30.2\%.
		
		\item [(d)] Evaluate the model fit for this regression. What does this tell us about the importance of yard signs versus other factors that are not modeled?
		
		The model has an $R^2 = 0.094$, meaning that the two covariates, whether a precinct was assigned lawn signs and whether it was adjacent to such precincts, explain only about 9.4\% of the variation in the vote share for Ken Cuccinelli.
		
		This suggests that yard signs have limited explanatory power, and that other potential factors influencing vote share are not included in this model.
		
		Thus, including additional relevant predictors would likely improve the explanatory power and overall fit of the model.
		
	\end{enumerate}
	
\end{document}